\section{Selected Projects}


\cventry
  {2019 --}
  {qubovert}
  {Python package (with C extension) for binary optimization}
  {}
  {}
  {
    \begin{itemize}
        \item Created \href{https://github.com/jtiosue/qubovert}{qubovert}, which is particularly designed to aid in converting optimization problems to a form that can be solved with quantum annealers and quantum optimization algorithms.
        \item qubovert can be installed with \texttt{pip install qubovert}, the source code is hosted at \url{github.com/jtiosue/qubovert}, and the documentation is hosted at \url{qubovert.readthedocs.io}.
        \item qubovert currently has \textbf{over 363k downloads} from \href{https://pypi.org/project/qubovert/}{PyPI}, 41 stars and 8 forks on GitHub, and has been \href{https://scholar.google.com/scholar?hl=en&as\_sdt=0\%2C33&q=qubovert&btnG=}{cited} in many research articles.
    \end{itemize}
  }


\cventry
  {2024}
  {(Deep) Q-learning in PyTorch}
  {}
  {}
  {}
  {
  \begin{itemize}
      \item \href{https://github.com/jtiosue/DotsBoxesQLearn}{PyTorch implementation} for self-playing reinforcement learning of a two player game.
  \end{itemize}
  }


\cventry
  {2023}
  {rcal}
  {Python package for review calibration}
  {}
  {}
  {
    \begin{itemize}
      \item Devised a novel review calibration algorithm (\href{https://github.com/jtiosue/rcal/blob/main/report/review_calibration.pdf}{written report})
      \item Implemented the algorithm in a Python \href{https://github.com/jtiosue/rcal}{package}.
    \end{itemize}
  }
    
    
\cventry
  {2019}
  {Powell bounded multivariate optimization}
  {SciPy contribution}
  {}
  {}
  {
    \begin{itemize}
      \item Authored \href{https://github.com/scipy/scipy/pull/10648}{pull request number 10648} on Python’s SciPy package. My contribution is included in the \href{https://github.com/scipy/scipy/releases/tag/v1.5.0}{1.5.0} release and later releases
      \item The pull request implements an additional feature for SciPy’s minimization method. I devised a bounded version of the standard unbounded Powell minimization method and found it to often perform much better than the other gradient-free minimizers. I then implemented this variant in SciPy’s software stack and created the pull request.
    \end{itemize}
  }
    
    
\cventry
  {2018}
  {Quantum Computer Simulator}
  {C\texttt{++} project}
  {}
  {}
  {
    \begin{itemize}
      \item Implemented a \href{https://github.com/jtiosue/Quantum-Computer-Simulator-with-Algorithms}{quantum computer simulator} in C\texttt{++}.
    \end{itemize}
  }



    